\clearpage
\begin{trapbox}

\begin{description}[style=nextline,leftmargin=2.8em,labelsep=0.5em,itemsep=0.45em]
  \item[\textbf{\textcolor{CTrap}{易错点一（指数大，值就一定大？）}}]

\begin{center}\begin{tikzpicture}[font=\small,>=stealth,line join=round]
\draw[gray!40,rounded corners=3pt] (0,0) rectangle (14.2,8.3);\node at (7.1,7.700000000000001) {例2：底数小于1，右边的点反而更低};
\begin{scope}[shift={(1.82,2.161111111111111)},x=2.4cm,y=2.5555555555555554cm]\draw[->,gray] (-0.3,0)--(2.2,0) node[right]{$x$};\draw[->,gray] (0,-0.2)--(0,1.6) node[above]{$y$};\draw[orange!80!black,thick,dashed,domain=-0.2:2,samples=180] plot(\x,{pow(.5,\x)});\fill[orange!85!black] (0.8,0.5743491774985174) circle(2.2pt);\fill[orange!85!black] (1.2,0.43527528164806206) circle(2.2pt);\node[below=5pt] at (0.8,0){$0.8$};\node[below=5pt] at (1.2,0){$1.2$};\draw[gray,densely dotted] (0.8,0)--(0.8,0.5743491774985174);\draw[gray,densely dotted] (1.2,0)--(1.2,0.43527528164806206);\end{scope}\node[anchor=west,align=left,text width=5.0cm] at (8.4,4.0) {图取$a=\frac12$\\[9pt]$0.8<1.2$\\[9pt]$a^{0.8}>a^{1.2}$\\[9pt]对任意$0<a<1$同样成立};
\node[text width=12.6cm,align=center] at (7.1,.65) {比较指数前先判断底数；下降的曲线把左右顺序变成相反的高低顺序。};
\end{tikzpicture}\end{center}

\textbf{观察：}左点高、右点低。\quad\textbf{原因：}乘0与1之间的正数会变小。

\textbf{考试使用：}先判底数，不等号与指数比较相反。


  \textbf{错误：}把 $(\frac12)^3$ 判得比 $(\frac12)^2$ 大。\textbf{纠正：}$(\frac12)^3=\frac18<\frac14=(\frac12)^2$。$0<a<1$ 时，每多乘一次小于 $1$ 的正数，值会变小；解同底不等式时必须先判断底数在哪一侧。

\end{description}
\end{trapbox}
\clearpage
\begin{trapbox}
\begin{description}[style=nextline,leftmargin=2.8em,labelsep=0.5em,itemsep=0.45em]
  \item[\textbf{\textcolor{CTrap}{易错点二（忽略底数范围）}}]

\begin{center}\begin{tikzpicture}[font=\small,>=stealth,line join=round]
\draw[gray!40,rounded corners=3pt] (0,0) rectangle (14.2,8.8);\node at (7.1,8.200000000000001) {底数范围有原因：a=1变成水平线，负底数不能全取实数};
\node at (3.1500000000000004,7.3) {$1^x=1$};\begin{scope}[shift={(2.8512195121951223,3.4981132075471697)},x=1.1951219512195124cm,y=0.6226415094339622cm]\draw[->,gray] (-1.8,0)--(2.3,0) node[right]{$x$};\draw[->,gray] (0,-0.8)--(0,4.5) node[above]{$y$};\draw[gray,thick,solid,domain=-1.5:2,samples=180] plot(\x,{1});\fill[orange!85!black] (0,1) circle(2pt);\fill[orange!85!black] (1,1) circle(2pt);\node[below left=5pt] at (0,0){$0$};\node[below=5pt] at (1,0){$1$};\end{scope}\node[align=center,text width=5.8cm] at (3.1500000000000004,1.7) {每处高度都为1\\单列为常数函数};\node at (9.95,7.3) {$(-2)^x$：只标整数示例};\begin{scope}[shift={(9.651219512195123,3.4981132075471697)},x=1.1951219512195124cm,y=0.6226415094339622cm]\draw[->,gray] (-1.8,0)--(2.3,0) node[right]{$x$};\draw[->,gray] (0,-0.8)--(0,4.5) node[above]{$y$};\fill[orange!85!black] (0,1) circle(2pt);\fill[orange!85!black] (2,4) circle(2pt);\node[below left=5pt] at (0,0){$0$};\node[below=5pt] at (0.5,0){$\frac12$};\node[below=5pt] at (2,0){$2$};\draw[gray,densely dotted] (.5,0)--(.5,3.5);\end{scope}\node[align=center,text width=5.8cm] at (9.95,1.7) {$x=\frac12$时无实数值\\不能连成全实数上的指数曲线};
\node[text width=12.6cm,align=center] at (7.1,.65) {两图同尺度；右图整数点只是示例，并不表示负底数幂的全部可用指数。};
\end{tikzpicture}\end{center}

\textbf{观察：}水平线无增减，负底数在某些实数处没有值。\quad\textbf{原因：}1的任何实数次方为1；负数不能在实数内开平方。

\textbf{考试使用：}先核对a>0且不等于1。


  \textbf{错误：}把 $1^x$ 或 $(-2)^x$ 直接当作定义域为 $\mathbb R$ 的指数函数。\textbf{纠正：}只讨论 $a>0$ 且 $a\neq1$。$1^x=1$ 是常数函数；负底数的实数次幂并非对所有实数 $x$ 都有定义。

\end{description}
\end{trapbox}
\clearpage
\begin{trapbox}
\begin{description}[style=nextline,leftmargin=2.8em,labelsep=0.5em,itemsep=0.45em]
  \item[\textbf{\textcolor{CTrap}{易错点三（把可变点当定点）}}]

\begin{center}\begin{tikzpicture}[font=\small,>=stealth,line join=round]
\draw[gray!40,rounded corners=3pt] (0,0) rectangle (14.2,8.3);\node at (7.1,7.700000000000001) {底数从4减到1/4：右侧降低，左侧升高，仍过(0,1)};
\begin{scope}[shift={(4.1,2.0951612903225807)},x=1.1538461538461537cm,y=0.7419354838709677cm]\draw[->,gray] (-2.6,0)--(2.6,0) node[right]{$x$};\draw[->,gray] (0,-0.6)--(0,5.6) node[above]{$y$};\draw[violet,thick,solid,domain=-2.3:1.131517202916897,samples=180] plot(\x,{pow(4,\x)});\draw[blue!75!black,thick,dashed,domain=-2.3:2.263034405833794,samples=180] plot(\x,{pow(2,\x)});\draw[orange!80!black,thick,densely dotted,domain=-2.263034405833794:2.3,samples=180] plot(\x,{pow(0.5,\x)});\draw[teal,thick,dashdotted,domain=-1.131517202916897:2.3,samples=180] plot(\x,{pow(0.25,\x)});\fill[orange!85!black] (0,1) circle(2.2pt);\draw[gray,densely dotted] (-2.3,1)--(2.3,1);\node[below=5pt] at (-2,0){$-2$};\node[below left=5pt] at (0,0){$0$};\node[below=5pt] at (2,0){$2$};\end{scope}\node[anchor=west,align=left,text width=5.0cm] at (8.4,4.0) {底数依次减小：\\[5pt]紫$4$、蓝$2$、橙$\frac12$、青$\frac14$\\[9pt]固定$x>0$：高度降低\\[9pt]固定$x<0$：高度升高\\[9pt]$x=0$：高度始终为$1$};
\node[text width=12.6cm,align=center] at (7.1,.65) {“顺时针”描述共同点附近的走势；曲线会改变弯曲程度，并非整条刚性旋转。};
\end{tikzpicture}\end{center}

\textbf{观察：}固定左右各一个横坐标，分别看高度。\quad\textbf{原因：}x>0时正幂保序；x<0时先取倒数，大小反向。

\textbf{考试使用：}同一侧比较曲线高低，判断底数大小；不能跨左右照搬。


  \textbf{错误：}声称每个 $y=a^x$ 都过 $(1,1)$ 或把 $(1,a)$ 叫作共同定点。\textbf{纠正：}$a^0=1$ 才说明\textbf{所有图像都经过 $(0,1)$}；当 $x=1$ 时，纵坐标 $a$ 会随底数变。

\end{description}
\end{trapbox}
\clearpage
\begin{trapbox}
\begin{description}[style=nextline,leftmargin=2.8em,labelsep=0.5em,itemsep=0.45em]
  \item[\textbf{\textcolor{CTrap}{易错点四（“接近 $0$”误当“等于 $0$”）}}]

\begin{center}\begin{tikzpicture}[font=\small,>=stealth,line join=round]
\draw[gray!40,rounded corners=3pt] (0,0) rectangle (14.2,8.3);\node at (7.1,7.700000000000001) {越往右越接近0，但每个点仍有正的高度};
\begin{scope}[shift={(1.6660377358490566,2.191176470588235)},x=1.1320754716981132cm,y=3.3823529411764706cm]\draw[->,gray] (-0.5,0)--(4.8,0) node[right]{$x$};\draw[->,gray] (0,-0.16)--(0,1.2) node[above]{$y$};\draw[orange!80!black,thick,dashed,domain=0:4.3,samples=180] plot(\x,{pow(2,-\x)});\fill[orange!85!black] (1,0.5) circle(2.2pt);\fill[orange!85!black] (3,0.125) circle(2.2pt);\fill[orange!85!black] (4,0.0625) circle(2.2pt);\node[below left=5pt] at (0,0){$0$};\node[below=5pt] at (1,0){$1$};\node[below=5pt] at (3,0){$3$};\node[below=5pt] at (4,0){$4$};\node[left=5pt] at (0,1){$1$};\draw[orange!80!black,densely dotted] (1,0)--(1,0.5);\draw[orange!80!black,densely dotted] (3,0)--(3,0.125);\draw[orange!80!black,densely dotted] (4,0)--(4,0.0625);\end{scope}\node[anchor=west,align=left,text width=5.0cm] at (8.4,4.0) {$y=(\frac12)^x$\\[9pt]$x=1:\ y=\frac12$\\[9pt]$x=3:\ y=\frac18$\\[9pt]$x=4:\ y=\frac1{16}>0$\\[9pt]值域$(0,+\infty)$，不含$0$};
\node[text width=12.6cm,align=center] at (7.1,.65) {每多乘一次二分之一，高度更小却仍为正；画在横轴附近，不表示落到了横轴上。};
\end{tikzpicture}\end{center}

\textbf{观察：}小高度仍在横轴上方。\quad\textbf{原因：}每次除以2，正数仍是正数。

\textbf{考试使用：}值域端点0用圆括号，不写最小值为0。


  \textbf{错误：}因为 $2^x$ 向左靠近 $0$，就说图像会穿过 $x$ 轴、值域包含 $0$。\textbf{纠正：}$a^x$ 对每个实数 $x$ 都大于 $0$；$0$ 是\textbf{靠近但取不到}的值，故值域是 $(0,+\infty)$ 而不是 $[0,+\infty)$。

\end{description}
\end{trapbox}
\clearpage
\begin{trapbox}
\begin{description}[style=nextline,leftmargin=2.8em,labelsep=0.5em,itemsep=0.45em]
  \item[\textbf{\textcolor{CTrap}{易错点五（关于 $y$ 轴对称的点写错）}}]

\begin{center}\begin{tikzpicture}[font=\small,>=stealth,line join=round]
\draw[gray!40,rounded corners=3pt] (0,0) rectangle (14.2,8.3);\node at (7.1,7.700000000000001) {关于y轴对称：横坐标变号，高度和距离都不变};
\begin{scope}[shift={(4.1,2.161111111111111)},x=1.0714285714285714cm,y=0.7301587301587301cm]\draw[->,gray] (-2.8,0)--(2.8,0) node[right]{$x$};\draw[->,gray] (0,-0.7)--(0,5.6) node[above]{$y$};\draw[blue!75!black,thick,solid,domain=-2.3:2.15,samples=180] plot(\x,{pow(2,\x)});\draw[orange!80!black,thick,dashed,domain=-2.3:2.15,samples=180] plot(\x,{pow(2,-\x)});\fill[orange!85!black] (-2,4) circle(2.2pt);\fill[orange!85!black] (2,4) circle(2.2pt);\fill[orange!85!black] (0,1) circle(2.2pt);\node[above=7pt] at (-1.25,4.65){$(-2,4)$};\node[above=7pt] at (1.25,4.65){$(2,4)$};\draw[<->,gray] (-2,4)--(0,4);\draw[<->,gray] (0,4)--(2,4);\node[below=4pt] at (-1,4){$2$};\node[below=4pt] at (1,4){$2$};\node[below=5pt] at (-2,0){$-2$};\node[below left=5pt] at (0,0){$0$};\node[below=5pt] at (2,0){$2$};\end{scope}\node[anchor=west,align=left,text width=5.0cm] at (8.4,4.0) {蓝：$f(x)=2^x$\\[8pt]橙：$g(x)=2^{-x}=f(-x)$\\[8pt]$(2,4)\leftrightarrow(-2,4)$\\[8pt]一般为$(t,f(t))$\\[4pt]$\leftrightarrow(-t,f(t))$};
\node[text width=12.6cm,align=center] at (7.1,.65) {同一水平线上的两点距y轴相等；从一条曲线可直接翻出另一条。};
\end{tikzpicture}\end{center}

\textbf{观察：}同一水平线的两点左右等距。\quad\textbf{原因：}倒数底数等于指数变号。

\textbf{考试使用：}对应点只改横坐标符号，交点为(0,1)。


  \textbf{错误：}把 $(t,f(t))$ 的对称点写成 $(-t,f(-t))$。\textbf{纠正：}关于 $y$ 轴对称时，纵坐标不能变，应为 $(-t,f(t))$。特别地，$(2,4)$ 在 $y=2^x$ 上，它在 $y=(\frac12)^x$ 上的对应点是 $(-2,4)$，\textbf{不是} $(-2,\frac14)$。

\end{description}
\end{trapbox}
\clearpage
\begin{trapbox}
\begin{description}[style=nextline,leftmargin=2.8em,labelsep=0.5em,itemsep=0.45em]
  \item[\textbf{\textcolor{CTrap}{易错点六（不同底数硬比较指数）}}]

\begin{center}\begin{tikzpicture}[font=\small,>=stealth,line join=round]
\draw[gray!40,rounded corners=3pt] (0,0) rectangle (14.2,8.3);\node at (7.1,7.700000000000001) {异底不能只比指数：2的三次方反而小于3的二次方};
\begin{scope}[shift={(2.6918367346938776,2.130597014925373)},x=1.2244897959183672cm,y=0.34328358208955223cm]\draw[->,gray] (-1.3,0)--(3.6,0) node[right]{$x$};\draw[->,gray] (0,-1.4)--(0,12) node[above]{$y$};\draw[blue!75!black,thick,solid,domain=-1:3.25,samples=180] plot(\x,{pow(2,\x)});\draw[violet,thick,dashed,domain=-1:2.25,samples=180] plot(\x,{pow(3,\x)});\fill[orange!85!black] (3,8) circle(2.2pt);\fill[orange!85!black] (2,9) circle(2.2pt);\node[below left=5pt] at (0,0){$0$};\node[below=5pt] at (2,0){$2$};\node[below=5pt] at (3,0){$3$};\draw[gray,densely dotted] (0,8)--(3,8);\draw[gray,densely dotted] (0,9)--(2,9);\end{scope}\node[anchor=west,align=left,text width=5.0cm] at (8.4,4.0) {蓝线上$(3,8)$：$2^3=8$\\[9pt]紫线上$(2,9)$：$3^2=9$\\[9pt]虽然$3>2$\\[9pt]却有$2^3<3^2$\\[9pt]先查是不是同底};
\node[text width=12.6cm,align=center] at (7.1,.65) {两个点来自不同的曲线；同一曲线的递增规律，不能直接拿来比较不同底数。};
\end{tikzpicture}\end{center}

\textbf{观察：}两个点在不同曲线上，高低不能由横坐标单独决定。\quad\textbf{原因：}只有同底时才应用同一函数的单调性。

\textbf{考试使用：}先化同底，或计算特殊幂、寻找中间值。


  \textbf{错误：}仅凭 $3>2$ 就断言 $2^3>3^2$。\textbf{纠正：}前者等于 $8$，后者等于 $9$；只有\textbf{底数相同且符合条件}时，才能直接靠指数大小判断。异底先化同底，或用 $1$、已知幂大小等作比较。

\end{description}
\end{trapbox}
\clearpage
\begin{trapbox}
\begin{description}[style=nextline,leftmargin=2.8em,labelsep=0.5em,itemsep=0.45em]
  \item[\textbf{\textcolor{CTrap}{易错点七（把平移后的定点、水平线照搬）}}]

\begin{center}\begin{tikzpicture}[font=\small,>=stealth,line join=round]
\draw[gray!40,rounded corners=3pt] (0,0) rectangle (14.2,8.3);\node at (7.1,7.700000000000001) {例6：共同点和平行靠近的水平线一起平移};
\begin{scope}[shift={(4.2875,2.079333333333333)},x=0.9375cm,y=0.6133333333333333cm]\draw[->,gray] (-3.4,0)--(3,0) node[right]{$x$};\draw[->,gray] (0,-0.7)--(0,6.8) node[above]{$y$};\draw[blue!75!black,thick,solid,domain=-3:1.6,samples=180] plot(\x,{pow(2,\x)});\draw[violet,thick,dashed,domain=-3:2.6,samples=180] plot(\x,{pow(2,\x-1)+3});\fill[orange!85!black] (0,1) circle(2.2pt);\fill[orange!85!black] (1,4) circle(2.2pt);\draw[violet,densely dotted] (-3.1,3)--(2.8,3);\node[above=6pt] at (-2.8,3){$3$};\draw[->,gray] (0,1)--(1,1)--(1,4);\node[below left=5pt] at (0,0){$0$};\node[below=5pt] at (1,0){$1$};\end{scope}\node[anchor=west,align=left,text width=5.0cm] at (8.4,4.0) {蓝：$y=2^x$\\[8pt]紫：$y=2^{x-1}+3$\\[8pt]$(0,1)\to(1,4)$\\[8pt]右移1，再上移3\\[8pt]$y=0\to y=3$\\[8pt]新值域$(3,+\infty)$};
\node[text width=12.6cm,align=center] at (7.1,.65) {每个点的横坐标加1、纵坐标加3；原来取不到的0也随之变成取不到的3。};
\end{tikzpicture}\end{center}

\textbf{观察：}点与水平线怎样一起移动。\quad\textbf{原因：}横坐标加1、纵坐标加3。

\textbf{考试使用：}先令指数为0找点，再由恒正确定值域。


  \textbf{错误：}认为 $y=2^{x-1}+3$ 仍过 $(0,1)$、仍以 $y=0$ 为渐近线。\textbf{纠正：}它过 $(1,4)$，靠近的是 $y=3$；\textbf{先看指数何时为 $0$，再看整体加了多少}。
\end{description}

\textbf{\textcolor{CTrap}{做题前的三项检查：}}底数是否合法？比较的是不是同底？翻转图像时是否只改变横坐标的符号？

\end{trapbox}
